\documentclass[10pt]{article}
\usepackage[T1]{fontenc}
\usepackage[utf8]{inputenc}
\usepackage[french]{babel}
\usepackage{amsmath,amssymb,mathrsfs,mathpazo}
\usepackage{geometry}
\geometry{a4paper,margin=2cm}
\usepackage{graphicx,array,tabularx,tabularray,colortbl,multirow,multicol,booktabs,enumitem,pifont,ulem,fancybox,diagbox,cancel}
\usepackage{pstricks,pstricks-add,pst-plot,pst-tree,pst-node,pst-text,pst-func,pst-math,pst-bspline,pst-3dplot,pst-eucl,pst-solides3d}
\let\pscalculate\undefined
\usepackage{pst-bezier}
\usepackage{tikz,pgfplots,tkz-tab}
\usetikzlibrary{arrows,arrows.meta,patterns,calc,positioning,trees,decorations.markings,shapes.geometric,shapes.misc,shapes.arrows,intersections,decorations.pathmorphing,decorations.pathreplacing,backgrounds,shadows,babel}
\pgfplotsset{compat=1.18}
\usepgfplotslibrary{fillbetween}
\usepackage[np]{numprint}
\usepackage{listings,makecell,pas-tableur,twoopt}
\let\barre\undefined
\usepackage[pstricks]{tablvar}
\definecolor{darkgreen}{RGB}{0,100,0}
\providecommand{\red}{\color{red}}
\providecommand{\blue}{\color{blue}}
\providecommand{\green}{\color{green}}
\providecommand{\black}{\color{black}}
\providecommand{\white}{\color{white}}
\providecommand{\gray}{\color{gray}}
\providecommand{\cyan}{\color{cyan}}
\providecommand{\magenta}{\color{magenta}}
\providecommand{\yellow}{\color{yellow}}
\providecommand{\violet}{\color{violet}}
\providecommand{\orange}{\color{orange}}
\providecommand{\euro}{\texteuro}
\providecommand{\EUR}{\texteuro}
\providecommand{\ieme}{\textsuperscript{e}}
\providecommand{\degres}{\ensuremath{^\circ}}
\providecommand{\degree}{\ensuremath{^\circ}}
\providecommand{\vv}[1]{\overrightarrow{#1}}
\providecommand{\Coord}[1]{#1}
\providecommand{\decosix}{}
\providecommand{\starredbullet}{\ensuremath{\bullet}}
\pagestyle{empty}
\setlength{\parindent}{0pt}
\usepackage{adjustbox}\DeclareUnicodeCharacter{2236}{:}\DeclareUnicodeCharacter{2190}{\ensuremath{\leftarrow}}
\providecommand{\R}{\mathbb{R}}
\providecommand{\N}{\mathbb{N}}
\providecommand{\Z}{\mathbb{Z}}
\providecommand{\Q}{\mathbb{Q}}
\providecommand{\D}{\mathbb{D}}
\providecommand{\C}{\mathbb{C}}
\providecommand{\e}{{\mathrm{e}}}
\providecommand{\vect}[1]{\overrightarrow{#1}}
\providecommand{\vectt}[1]{\overrightarrow{\mathrm{#1}}}
\providecommand{\vvt}[1]{\overrightarrow{\mathrm{#1}}}
\providecommand{\np}[2][]{{#2~#1}}
\providecommand{\vv}[1]{{\overrightarrow{#1}}}
\renewcommand{\d}{{\mathrm{d}}}
\providecommand{\degres}{{^\circ}}
\providecommand{\degre}{{^\circ}}
\providecommand{\blue}{\color{blue}}
\providecommand{\red}{\color{red}}
\providecommand{\enskip}{\quad}
\providecommand{\llbracket}{\lbrack\!\lbrack}
\providecommand{\rrbracket}{\rbrack\!\rbrack}
\providecommand{\euro}{{\text{€}}}
\providecommand{\no}{{\text{n°}}}
\providecommand{\interval}[2]{{[#1;#2]}}
\providecommand{\intervFO}[2]{{[#1;#2[}}
\providecommand{\dotfill}{\ldots}

\providecommand{\ds}{\displaystyle}
\providecommand{\textding}[1]{#1}
\providecommand{\Oij}{\(\mathrm{O};\vec{\imath},\vec{\jmath}\)}
\providecommand{\Oijk}{\(\mathrm{O};\vec{\imath},\vec{\jmath},\vec{k}\)}
\providecommand{\Ouv}{\(\mathrm{O};\vec{u},\vec{v}\)}
\providecommand{\intervFF}[2]{\ensuremath{[#1;#2]}}
\renewcommand{\euro}{\ifmmode\text{\texteuro}\else\texteuro\fi}
\DeclareUnicodeCharacter{00B0}{\ifmmode^{\circ}\else\textdegree\fi}
\DeclareUnicodeCharacter{20AC}{\ifmmode\text{\texteuro}\else\texteuro\fi}
% Énoncés : Éducation nationale. Transcription et corrigés : APMEP, Denis Vergès et contributeurs. Découpage et classement : CoopMaths, Sébastien Lozano, Éric Elter et contributeurs. Adaptation de la présentation : Maths Explorer.
% CC BY-SA : https://coopmaths.fr/www/about/
% https://coopmaths.fr/alea/static/bac/2026/tex/bac_2026_06_sujet2_metropole_3.tex
\begin{document}
\noindent\textit{APMEP / CoopMaths. CC BY-SA. Présentation adaptée par Maths Explorer.}\par\medskip
\emph{Pour chacune des affirmations suivantes, indiquer si elle est vraie ou fausse.\\
	Justifier chaque réponse.\\
	Une réponse non justifiée ne rapporte aucun point.}

\medskip

	On désigne par \og musiciens professionnels \fg{} l'ensemble des femmes et hommes pratiquant un instrument de manière professionnelle.

	\begin{enumerate}[itemsep=5mm]
		\item Une étude a été réalisée auprès des musiciens professionnels pratiquant différents genres musicaux (Pop, Rap, Rock, Classique, etc.).


		\begin{list}{\textbullet}{Cette étude indique que:}
		\item $52\,\%$ des musiciens professionnels jouent de la Pop ;
		\item parmi les musiciens professionnels jouant de la Pop, $32\,\%$ sont des femmes ;
		\item $20\,\%$ des musiciens professionnels sont des femmes.
			\end{list}

		On choisit au hasard une personne dans l'ensemble des musiciens professionnels.

		\begin{list}{\textbullet}{On note:}
			\item $O$ l'évènement : \og la personne joue de la Pop \fg{} ;
			\item $F$ l'évènement : \og la personne est une femme \fg{}.
		\end{list}

		\textbf{Affirmation 1 :} Si la personne choisie est une femme, alors la probabilité qu'elle joue de la Pop est de 0,832.

		\item La proportion de musiciens professionnels au sein de la population active dans le secteur de la culture est estimée à $6,2\,\%$.

		On prélève au hasard et de manière indépendante un échantillon de \np{5000} personnes au sein de cette population.

		La population considérée est assez importante pour qu'on puisse assimiler ce prélèvement à un tirage avec remise.

		On considère la variable aléatoire $X$ qui, à tout échantillon ainsi défini, associe le nombre de musiciens professionnels dans l'échantillon.

		\medskip

		\textbf{Affirmation 2 :} En arrondissant le résultat au dixième, on a \quad $P(X \leqslant  340) \approx 0,4$.

		\medskip

		\textbf{Affirmation 3 :} En appliquant l'inégalité de Bienaymé-Tchebychev, on peut affirmer qu'il y a plus de $95\,\%$ de chance que, sur cet échantillon de \np{5000} actifs dans le secteur de la culture, le nombre de musiciens professionnels soit strictement compris entre 230 et 390.

		\item Un jeu sur le thème de la culture de la Pop est proposé à 10 personnes dont 4 sont des musiciens professionnels.

		Deux équipes de 5 joueurs s'affrontent.

		Elles sont chacune constituées de 2 musiciens professionnels et de 3 personnes qui ne le sont pas.

		\textbf{Affirmation 4 :} On peut former ainsi 120 équipes différentes.
	\end{enumerate}

\end{document}
